\documentclass[a4paper,12pt]{report}
\usepackage[utf8]{inputenc}
\usepackage[italian]{babel}
\usepackage{graphicx}
\usepackage[table]{xcolor} % AGGIUNTO [table] per \rowcolor
\usepackage{tikz}
\usepackage{amsmath, amssymb}
\usepackage{listings}
\usepackage{hyperref}
\usepackage{geometry}
\usepackage{tcolorbox}
\usepackage{tabularx}
\usepackage{pgfplots}
\pgfplotsset{compat=1.18}

% Librerie TikZ necessarie
\usetikzlibrary{shapes, arrows.meta, positioning, calc, shadows, fit}

% --- DEFINIZIONE LINGUAGGIO KOTLIN PER LISTINGS ---
\lstdefinelanguage{Kotlin}{
    keywords={package, as, typealias, this, super, val, var, fun, for, null, true, false, is, in, throw, return, break, continue, object, if, else, while, do, try, when, class, interface, enum, self, companion, override, private, public, internal, protected, import, lateinit, by},
    keywordstyle=\color{blue}\bfseries,
    ndkeywords={@Entity, @Dao, @Query, @Insert, @Update, @Delete, @Database, @PrimaryKey, @Composable, @Volatile, @remember, @mutableStateOf},
    ndkeywordstyle=\color{orange}\bfseries,
    identifierstyle=\color{black},
    sensitive=true,
    comment=[l]{//},
    morecomment=[s]{/*}{*/},
    commentstyle=\color{gray}\ttfamily,
    stringstyle=\color{green!50!black}\ttfamily,
    morestring=[b]",
    morestring=[s]{"""}{"""}
}

\lstset{
    language=Kotlin, % Impostato come default
    backgroundcolor=\color{gray!5},
    basicstyle=\ttfamily\footnotesize,
    breaklines=true,
    numbers=left,
    numberstyle=\tiny\color{gray},
    frame=single,
    tabsize=2
}

\hypersetup{
    colorlinks=true,
    linkcolor=black,
    filecolor=black,
    urlcolor=blue,
    citecolor=black,
}

\title{\textbf{Appunti di Programmazione di Dispositivi Mobili}}
\author{G.C.\\ Università degli Studi dell'Insubria}
\date{Anno Accademico 2025-2026}

\begin{document}

    \maketitle

    \begin{abstract}
        Appunti completi per l'esame di PDM basati sulle slide del Prof. Gallo.
    \end{abstract}

    \tableofcontents
    \newpage

    \chapter{Introduzione allo Sviluppo Mobile e Android}

\section{Il mercato Mobile e lo Sviluppo Cross-Platform}
Oggi il mercato dei sistemi operativi mobili è dominato da due attori principali: \textbf{Android} (Google) che detiene circa il 70-74\% del mercato globale, e \textbf{iOS} (Apple) con il restante 25-30\%.

Quando si decide di sviluppare un'applicazione mobile, ci si trova davanti a un bivio architetturale:
\begin{itemize}
    \item \textbf{Sviluppo Nativo:} Scrivere codice specifico per ogni piattaforma (Kotlin/Java per Android, Swift/Objective-C per iOS).
    Garantisce le massime performance e accesso immediato alle API del dispositivo.
    \item \textbf{Sviluppo Cross-Platform:} Scrivere il codice una sola volta e deployarlo su più piattaforme.
    Abbassa i costi e i tempi di sviluppo, ma può portare a compromessi sulle prestazioni e ritardi nel supporto alle nuove API\@.
\end{itemize}

I framework Cross-Platform più popolari ad oggi includono:
\begin{itemize}
    \item \textbf{React Native (Facebook):} Basato su JavaScript.
    \item \textbf{Flutter (Google):} Basato sul linguaggio Dart, rendering custom (non usa le UI native ma le ridisegna).
    \item \textbf{Kotlin Multiplatform Mobile (KMM):} Creato da JetBrains.
    Non è un framework UI, ma un \textit{SDK} che permette di condividere la logica di business (networking, database, logica)
    in Kotlin, mantenendo lo sviluppo della UI strettamente nativo per Android e iOS.
\end{itemize}

\section{Architettura di un'App Android e Componenti Core}
Un'applicazione Android non ha un singolo ``entry point'' (come il classico \texttt{main()} in C o Java desktop), ma è composta da quattro blocchi fondamentali (Componenti):
\begin{enumerate}
    \item \textbf{Activities:} Rappresentano le singole schermate con un'interfaccia utente (UI).
    \item \textbf{Services:} Eseguono operazioni di lunga durata in background, senza UI (es.
    riproduzione musicale, download).
    \item \textbf{Broadcast Receivers:} Ascoltano e reagiscono a messaggi o eventi di sistema (es.
    "batteria scarica", "schermo sbloccato").
    \item \textbf{Content Providers:} Gestiscono e condividono un set di dati tra diverse applicazioni in modo sicuro (es.
    la rubrica dei contatti).
\end{enumerate}

Lo sviluppo moderno (Modern Android Development) incoraggia l'uso di librerie \textbf{Jetpack} (es. ViewModel per lo stato, WorkManager per il background, Room per i DB) che semplificano enormemente la gestione del ciclo di vita di questi componenti.

\section{Android Studio e il processo di Build (Gradle)}
\textbf{Android Studio} è l'IDE ufficiale di Google basato su IntelliJ IDEA. Dal 2015 ha sostituito Eclipse.

\subsection{Gradle e Architettura Hardware}
Nei quiz è importante ricordare che:
\begin{itemize}
    \item \textbf{Gradle} è un sistema di build automation open-source.
    Costruito sui concetti di Apache Ant e Maven, utilizza un \textbf{DSL (Domain-Specific Language)} basato su Groovy o Kotlin.
    \item I dispositivi mobili utilizzano spesso l'\textbf{architettura ARM big.LITTLE}: un'architettura che combina
    nuclei di processori di dimensioni diverse per \textbf{ottimizzare le prestazioni e l'efficienza energetica}.
\end{itemize}

\subsection{Dall'Idea all'Esecuzione: APK, DEX e ART}
Nei quiz è fondamentale conoscere come il codice diventa un'app eseguibile:
\begin{enumerate}
    \item Il codice sorgente (Kotlin/Java) viene compilato nel bytecode standard di Java (\texttt{.class}).
    \item Il compilatore Android converte i file \texttt{.class} nel formato \textbf{DEX (Dalvik Executable)}, ottimizzato per dispositivi mobili.
    \item Il tool \texttt{aapt} unisce i file DEX e le risorse compilate per creare l'archivio \textbf{APK} (Android Application Package).
\end{enumerate}

\begin{tcolorbox}[colback=blue!5!white,colframe=blue!75!black,title=Android Runtime (ART) vs Dalvik]
    In passato, Android usava la macchina virtuale \textbf{Dalvik} (compilazione JIT - Just In Time).
    Oggi utilizza \textbf{ART (Android Runtime)}.
    ART compila l'app in codice macchina nativo al momento dell'installazione tramite la tecnica \textbf{AOT (Ahead-Of-Time) compilation}.
    \newline\textit{Vantaggi:} Esecuzione molto più veloce. \newline\textit{Svantaggi:} Tempi di installazione leggermente più lunghi e maggiore spazio occupato sul disco.
\end{tcolorbox}

\section{Figma: UI vs UX e Prototipazione}
Prima di scrivere il codice, la fase di design è fondamentale.
\begin{itemize}
    \item \textbf{UX (User Experience):} Struttura logica, gerarchie, posizionamento e flussi funzionali.
    Analizza i comportamenti reali degli utenti.
    \item \textbf{UI (User Interface):} Struttura visiva, colori, font, icone.
    Assicura che il prodotto sia visivamente stimolante.
\end{itemize}

\begin{tcolorbox}[colback=yellow!10!white,colframe=yellow!70!black,title=Figma nei Quiz]
    \textbf{Figma} è una piattaforma di design collaborativo \textbf{basata su cloud} (Web-based tramite WebGL/WASM).
    Utilizza i \textbf{CRDT} (Conflict-free Replicated Data Types) e \textbf{WebSockets} per permettere a più persone
    di lavorare sullo stesso documento in \textit{real-time} senza conflitti di file (approccio \textit{Last-Writer-Wins}).
    Consente di creare \textbf{veri prototipi} per testare l'esperienza utente prima di scrivere codice.
    \begin{itemize}
        \item \textbf{Caratteristica principale:} È \textit{Web-based} (WASM + WebGL).
        A differenza di software storici come Sketch, non richiede installazioni locali o file \texttt{.sketch} sul disco.
        \item Il motore di rendering sotto il cofano è scritto in \textbf{C++} e transpilato in WebAssembly, permettendo la
        gestione di grafica vettoriale con la potenza di un videogioco direttamente nel browser.
    \end{itemize}
\end{tcolorbox}

\subsection{Il contenuto di un APK}
Da un punto di vista dei file, un APK è un archivio ZIP che contiene:
\begin{itemize}
    \item \texttt{classes.dex}: Il codice eseguibile dell'applicazione.
    \item \texttt{res/}: Le risorse non compilate.
    \item \texttt{resources.arsc}: Le risorse compilate (tabelle di lookup generate da \texttt{aapt}).
    \item \texttt{META-INF/}: Informazioni sulle firme digitali.
    \item \textbf{\texttt{AndroidManifest.xml}}: Il file fondamentale che descrive l'app al sistema operativo Android.
\end{itemize}

\section{AndroidManifest.xml}
Ogni app Android deve avere un file chiamato \texttt{AndroidManifest.xml} nella root del progetto. 
Definisce le caratteristiche essenziali dell'applicazione necessarie ad Android prima ancora che l'app venga avviata:
\begin{itemize}
    \item Dichiara il \textit{package name} (che funge da ID univoco dell'app).
    \item \textbf{Dichiara i Componenti:} Tutte le \texttt{<activity>}, \texttt{<service>}, \texttt{<receiver>} e \texttt{<provider>} devono essere dichiarati qui. Se crei una classe Activity in Kotlin ma non la dichiari nel Manifest, l'app andrà in crash se cerchi di avviarla.
    \item Dichiara i \textbf{permessi} (\texttt{<uses-permission>}) richiesti (es. INTERNET, fotocamera).
    \item Dichiara le caratteristiche hardware richieste (es. touch screen, NFC).
\end{itemize}

Un costrutto chiave all'interno del manifest è l'\textbf{Intent-Filter}. Per indicare al sistema operativo qual è la prima Activity da lanciare all'avvio dell'app (l'Activity \textit{Home}), si usa la seguente sintassi:
\begin{lstlisting}[language=XML]
<activity android:name=".MainActivity" android:exported="true">
    <intent-filter>
        <action android:name="android.intent.action.MAIN" />
        <category android:name="android.intent.category.LAUNCHER" />
    </intent-filter>
</activity>
\end{lstlisting}
Il flag \texttt{android:exported="true"} indica che questa Activity può essere avviata da app esterne (incluso il launcher (la schermata home) del sistema operativo).
    \chapter{Il linguaggio Kotlin}

\section{Variabili e Tipi Base}
In Kotlin \textbf{tutto è un oggetto} (non esistono tipi primitivi come in Java). Il compilatore non effettua conversioni implicite (es. non puoi assegnare un \texttt{Int} a un \texttt{Long} senza \texttt{.toLong()}).

\begin{itemize}
    \item \textbf{\texttt{val}:} Immutabile (sola lettura).
    \item \textbf{\texttt{var}:} Mutabile.
\end{itemize}

\section{Null Safety}
Per evitare il crash dell'app (\texttt{NullPointerException}), Kotlin introduce tipi nullable.

\begin{center}
    \begin{tikzpicture}[
        node distance=1cm and 1.5cm,
        box/.style={rectangle, draw, rounded corners, fill=white, text width=3.5cm, align=center, minimum height=1.2cm, drop shadow},
        arrow/.style={-Stealth, thick}
    ]

% Nodes
        \node[box, fill=blue!10] (nullable) {\textbf{String?}\\ (Può essere null)};
        \node[box, below=1.5cm of nullable] (safecall) {\textbf{Safe Call (?.)}\\Ritorna null se la variabile è null};
        \node[box, left=of safecall] (elvis) {\textbf{Elvis (?:)}\\Valore di default};
        \node[box, right=of safecall] (bang) {\textbf{Assertion (!!)}\\Forza: rischio crash!};

% Arrows
        \draw[arrow] (nullable) -- (safecall);
        \draw[arrow] (nullable) -| (elvis);
        \draw[arrow] (nullable) -| (bang);
    \end{tikzpicture}
\end{center}

\section{Classi e Costruttori}
L'ordine di esecuzione è un tipico trabocchetto d'esame:
\begin{enumerate}
    \item Inizializzazione delle proprietà della classe.
    \item Blocchi \textbf{\texttt{init}} (nell'ordine in cui appaiono).
    \item \textbf{Costruttore secondario} (se invocato).
\end{enumerate}

\begin{tcolorbox}[colback=red!5!white,colframe=red!75!black,title=Regola d'oro]
    Il costruttore secondario \textbf{deve sempre delegare} al primario usando \texttt{this()}.
\end{tcolorbox}

\section{Interfacce vs Classi Astratte}
\begin{table}[h]
    \centering
    \begin{tabularx}{\textwidth}{|l|X|X|}
\hline
\rowcolor{gray!20} \textbf{Caratteristica} & \textbf{Interfaccia} & \textbf{Classe Astratta} \\ \hline
\textbf{Stato} & No (no variabili istanza) & Sì (variabili istanza) \\ \hline
\textbf{Ereditarietà} & Molteplice & Singola \\ \hline
\textbf{Metodi} & Astratti o con default & Astratti o concreti \\ \hline
    \end{tabularx}
\end{table}

\section{Functional Programming}
Kotlin permette di trattare le funzioni come variabili.
\begin{itemize}
    \item \textbf{Lambda:} Funzioni anonime \texttt{\{ a, b -> a + b \}}.
    \item \textbf{Uso di \texttt{it}:} Se la lambda ha un solo parametro, puoi omettere il nome e usare la keyword \texttt{it}.
\end{itemize}

\begin{lstlisting}[caption={Esempio di Lambda e Filter}]
val numeri = listOf(3, 8, 12, 5)
// Versione estesa
val maggiori7 = numeri.filter { n -> n > 7 }
// Versione compatta con 'it'
val maggiori7Compatta = numeri.filter { it > 7 }
\end{lstlisting}
    \chapter{Gestione delle Risorse e UI}

Un'applicazione Android è formata da Codice (Java/Kotlin) e \textbf{Risorse}.
Le risorse sono tutto ciò che non è codice logico: file XML di layout, pacchetti lingua, immagini, file audio, stili e colori.
L'obiettivo è separare la presentazione (Layout) dalla logica (Management).

\section{La cartella \texttt{res/} e il file \texttt{R.java}}
Le risorse vengono salvate nella cartella \texttt{res/} e suddivise in sottocartelle specifiche:
\begin{itemize}
    \item \texttt{layout/}: file XML che definiscono le interfacce grafiche.
    \item \texttt{values/}: contiene file XML che memorizzano valori semplici come stringhe (\texttt{strings.xml}), colori (\texttt{colors.xml}), dimensioni (\texttt{dimens.xml}) e stili (\texttt{styles.xml}).
    \item \texttt{drawable/}: per immagini (.png, .jpg), grafiche vettoriali o shape XML.
    \item \texttt{mipmap/}: dedicato esclusivamente alle icone dell'applicazione (launcher icons).
\end{itemize}

\subsection{Meccanismo di Risoluzione (R.java)}
Al momento della build, il tool \texttt{aapt} assegna ad ogni risorsa un ID intero univoco in formato esadecimale (es. \texttt{0x7f05000b}) e lo memorizza in una classe autogenerata chiamata \textbf{\texttt{R.java}}.
\begin{itemize}
    \item Da XML ci si riferisce alle risorse con la sintassi: \texttt{@tipo/nome} (es. \texttt{@string/app\_name}).
    \item Da codice Kotlin/Java ci si riferisce con: \texttt{R.tipo.nome} (es. \texttt{R.id.button1}).
\end{itemize}

\section{Risoluzione, Densità e le Unità DP e SP}
Android gira su migliaia di dispositivi diversi.
L'applicazione deve mantenere l'indipendenza dalla densità dello schermo ("Density Independence").

\begin{tcolorbox}[colback=green!5!white,colframe=green!50!black,title=Definizioni Chiave per i Quiz]
    \textbf{Resolution (Risoluzione):} Numero totale di pixel dello schermo (es.
    1920x1080).
    Due schermi delle stesse dimensioni fisiche possono avere risoluzioni diverse.\\
    \textbf{Density (DPI / PPI):} Punti per pollice.
    Misura quanto i pixel sono vicini fisicamente l'uno all'altro.
\end{tcolorbox}

Se creassimo un pulsante largo "100 pixel" (\texttt{100px}), questo apparirebbe grande su un vecchio dispositivo e minuscolo su uno smartphone moderno ad altissima definizione.
Per questo si usano i \textbf{Density Independent Pixels (DP)}.

\subsection{La Formula dei DP}
La baseline di Android è fissata a una densità media (\textbf{mdpi} = 160 dpi).
In questa specifica densità, \textbf{1 dp = 1 px}.
La formula per convertire da DP a Pixel hardware è:
\begin{equation}
    px = dp \times \left( \frac{dpi}{160} \right)
\end{equation}
Esempio: su uno schermo \texttt{xxhdpi} (480 dpi), un elemento di 10 dp corrisponderà a $10 \times (480/160) = 30px$ fisici.

Per i testi si utilizza invece \textbf{SP (Scale-independent Pixels)}, che si comporta matematicamente come il DP, ma
scala dinamicamente in base alle preferenze di dimensione del font impostate dall'utente nel sistema operativo.

\begin{tcolorbox}[colback=white,colframe=black,title=Calcolo della dimensione fisica (Quiz)]
    Nei quiz viene chiesto di calcolare la dimensione fisica in pollici (\textit{inches}) di un elemento. La formula è:
    \begin{center}
        $\text{Physical Size (in)} = \frac{\text{Pixel Reali}}{\text{DPI del dispositivo}}$
    \end{center}
    Ricordando che $\text{Pixel Reali} = \text{DP} \times (\frac{\text{DPI}}{160})$.\\
    \textbf{Esempio:} 100 dp su uno schermo a 240 dpi. \\
    1) Pixel = $100 \times (240/160) = 150px$. \\
    2) Physical Size = $150 / 240 = 0.625$ pollici.
\end{tcolorbox}

\section{Qualificatori e Risorse Alternative}
Per gestire dispositivi diversi (es.
orientamento, lingua), si creano cartelle "alternative" specificando un \textbf{qualificatore} nel nome della cartella: \texttt{<nome\_risorsa>-<qualificatore>}.
Esempi:
\begin{itemize}
    \item \texttt{res/layout/} (Layout di default per smartphone portrait)
    \item \texttt{res/layout-land/} (Layout che Android carica automaticamente se si ruota il telefono)
    \item \texttt{res/values-en/} (Stringhe in lingua inglese)
\end{itemize}
A runtime, Android cerca il "best match".
Se un qualificatore richiesto contraddice la configurazione del dispositivo (es. cerco layout per \texttt{fr} ma l'app gira in \texttt{en-GB}), Android elimina quelle directory e usa i fallback di default.

\section{Stili (Styles) vs Temi (Themes)}
Entrambi si definiscono tramite i tag \texttt{<style>} dentro \texttt{res/values/styles.xml}.
Possono ereditare proprietà da un genitore tramite l'attributo \texttt{parent}.
La differenza sta nell'\textbf{applicazione}:
\begin{itemize}
    \item \textbf{Style:} Viene applicato a una singola \texttt{View} (es. un bottone) tramite l'attributo \texttt{style="@style/MioStile"}.
    \item \textbf{Theme:} È un'estensione dello style, ma viene applicato all'intera applicazione o a una singola
    Activity nel file \texttt{AndroidManifest.xml} (es. \texttt{android:theme="@style/MioTema"}).
    Un tema modifica parametri globali come la barra di navigazione, lo status bar o i colori primari.
\end{itemize}

\vspace{0.5cm}
\begin{center}
    \begin{tikzpicture}
        % UI vs UX Box
        \node[draw, rounded corners, fill=blue!5, text width=6cm, align=center] (UX) at (-3.5,0) {
            \textbf{UX (User Experience)}\\
            Si occupa del vissuto dell'utente.
            Logica, gerarchia, fluidità dei processi, wireframe.
        };

        \node[draw, rounded corners, fill=orange!5, text width=6cm, align=center] (UI) at (3.5,0) {
            \textbf{UI (User Interface)}\\
            Il layout grafico visivo.
            Colori, font, bordi, immagini e icone.
        };

        \node[draw, fill=gray!20, text width=8cm, align=center] (Figma) at (0,-2.5) {
            \textbf{Figma}\\
            Piattaforma WebGL/WASM cloud-based per il design collaborativo.
            Usa un approccio \textit{Multiplayer Real-Time} basato su CRDT per prevenire conflitti.
        };

        \draw[->, thick, dashed] (UX) -- (Figma);
        \draw[->, thick, dashed] (UI) -- (Figma);
    \end{tikzpicture}
\end{center}
    % --- START OF FILE 04_activity_intents.tex ---

\chapter{Activity, Intent e Ciclo di Vita}

L'\textbf{Activity} è uno dei mattoni fondamentali (\textit{building blocks}) delle app Android. Rappresenta una singola schermata con un'interfaccia utente. Funge da punto di ingresso (\textit{entry point}) per l'interazione dell'utente con l'app.

\section{Gli Intent}
Un \textbf{Intent} è un oggetto di messaggistica (\textit{message-passing framework}) che permette a componenti diversi di richiedere azioni ad altri componenti, sia all'interno della stessa app che verso app esterne.

Possono essere di due tipi principali:
\begin{itemize}
    \item \textbf{Intent Espliciti (Explicit Intents):} Specificano \textit{esattamente} la classe del componente da avviare (es. \texttt{ActivityB::class.java}). Si usano quasi sempre all'interno della propria app.
    \item \textbf{Intent Impliciti (Implicit Intents):} Non specificano il componente esatto, ma indicano un'\textbf{azione} generica (es. \texttt{ACTION\_VIEW}) e dei dati (es. un URL). Il sistema Android verifica gli \textit{Intent-Filter} delle app installate e sceglie quella adatta (o chiede all'utente quale usare).
\end{itemize}

\begin{lstlisting}[caption={Esempi di Intent}, label={lst:intents}]
// 1. Intent Esplicito
val intentEsplicito = Intent(this, SecondActivity::class.java)
startActivity(intentEsplicito)

// 2. Intent Implicito (aprire una pagina web)
val uri = Uri.parse("http://www.uninsubria.it")
val intentImplicito = Intent(Intent.ACTION_VIEW, uri)
startActivity(intentImplicito)
\end{lstlisting}

\section{Activity Result API (Risultati dalle Activity)}
Spesso un'Activity ne avvia un'altra per ottenere un risultato (es. scattare una foto, scegliere un contatto).

\begin{tcolorbox}[colback=red!5!white,colframe=red!75!black,title=Attenzione all'esame: Deprecazione!]
    Il vecchio metodo \texttt{startActivityForResult()} e il relativo callback \texttt{onActivityResult()} sono \textbf{deprecati}. Erano problematici perché, se l'app veniva messa in background e uccisa dal sistema, il ritorno del risultato si perdeva.
\end{tcolorbox}

Il nuovo standard è l'\textbf{Activity Result API} (da \texttt{androidx 1.2.0}), che si divide in 3 passaggi fondamentali:
\begin{enumerate}
    \item \textbf{Contract:} Definisce input e output (es. \texttt{StartActivityForResult()}).
    \item \textbf{Register:} Si registra il contratto nella fase di inizializzazione chiamando \texttt{registerForActivityResult()}, che restituisce un \texttt{ActivityResultLauncher}.
    \item \textbf{Launch:} Si invoca il metodo \texttt{launch()} sull'oggetto Launcher quando si vuole avviare l'azione.
\end{enumerate}

\begin{lstlisting}[caption={Ottenere un risultato con Activity Result API (Esatto da Quiz)}]
// 1. Registrazione del contratto (da fare come attributo della classe)
val launcher = registerForActivityResult(
    ActivityResultContracts.StartActivityForResult()
) { result ->
    // Questo e il blocco dell'ActivityResultCallback
    if (result.resultCode == Activity.RESULT_OK) {
        val data = result.data?.getStringExtra("RISPOSTA")
    }
}

// 2. Lancio dell'Intent (es. dentro un listener)
button.setOnClickListener {
    val intent = Intent(this, SecondActivity::class.java)
    launcher.launch(intent) // L'istruzione per far partire la richiesta
}
\end{lstlisting}
\begin{tcolorbox}[colback=yellow!10!white,colframe=yellow!70!black,title=Perché \texttt{startActivityForResult} è deprecato?]
    Il vecchio metodo è stato deprecato perché il codice era strettamente accoppiato al ciclo di vita dell'Activity, generando bug (es. crash o perdita di dati) se l'Activity chiamante veniva distrutta dal sistema mentre la seconda Activity era aperta.
\end{tcolorbox}

\section{Il Ciclo di Vita dell'Activity (Lifecycle)}
Comprendere il ciclo di vita è essenziale. Android notifica i cambiamenti di stato dell'Activity invocando metodi di \textit{callback}.

\begin{center}
    \begin{tikzpicture}[
        node distance=0.8cm and 1cm,
        state/.style={rectangle, draw, rounded corners, fill=white, text width=3.5cm, align=center, drop shadow},
        arrow/.style={-Stealth, thick}
    ]
        \node[state, fill=blue!10] (create) {\textbf{onCreate()}\\Setup iniziale (UI)};
        \node[state, fill=cyan!10, below=of create] (start) {\textbf{onStart()}\\Activity Visibile};
        \node[state, fill=green!20, below=of start] (resume) {\textbf{onResume()}\\Focus / Interattiva};
        \node[state, fill=yellow!20, below=of resume] (pause) {\textbf{onPause()}\\Parz. visibile / No Focus};
        \node[state, fill=orange!20, below=of pause] (stop) {\textbf{onStop()}\\Completamente nascosta};
        \node[state, fill=red!20, below=of stop] (destroy) {\textbf{onDestroy()}\\Distruzione};

        \draw[arrow] (create) -- (start);
        \draw[arrow] (start) -- (resume);
        \draw[arrow] (resume) -- (pause);
        \draw[arrow] (pause) -- (stop);
        \draw[arrow] (stop) -- (destroy);

        \draw[arrow, dashed] (pause.east) -- ++(1.5,0) |- node[right, near start] {Ritorno in Foreground} (resume.east);
        \draw[arrow, dashed] (stop.west) -- ++(-1.5,0) |- node[left, near start] {Riavvio} (start.west);

        % Box laterali per le vite
        \draw[decorate, decoration={brace, amplitude=10pt}, thick]
        (resume.north east) -- (pause.south east) node[midway, right=15pt, align=left] {\textbf{Foreground Lifetime}\\(Multitasking, Popups)};

        \draw[decorate, decoration={brace, amplitude=10pt, mirror}, thick]
        (start.north west) -- (stop.south west) node[midway, left=15pt, align=right] {\textbf{Visible Lifetime}\\(L'utente vede l'app)};
    \end{tikzpicture}
\end{center}

Dettagli chiave sui metodi:
\begin{itemize}
    \item \textbf{\texttt{onCreate()}}: Chiamato una sola volta per instanziare l'UI (\texttt{setContentView}).
    \item \textbf{\texttt{onResume()} / \texttt{onPause()}}: Segnano i confini del \textit{Foreground lifetime}. In \texttt{onPause} l'app è ancora visibile ma non ha il focus (es. in modalità Multi-finestra, o se c'è un Dialog semi-trasparente sopra).
    \item \textbf{\texttt{onStop()}}: L'app non è più visibile (es. utente preme il tasto Home o apre un'altra app a schermo intero).
\end{itemize}

\section{Gestione delle Configuratrion Changes}
Quando avviene un cambiamento di configurazione a runtime (es. \textbf{rotazione dello schermo}, cambio della lingua, cambio modalità dark/light), il comportamento di default di Android è \textbf{distruggere e ricreare l'Activity} (\texttt{onPause} $\rightarrow$ \texttt{onStop} $\rightarrow$ \texttt{onDestroy} $\rightarrow$ \texttt{onCreate}).

\textbf{Il Problema:} Questo fa perdere tutte le variabili locali e lo stato non salvato.

Le soluzioni in ordine di correttezza (architetturale):
\begin{enumerate}
    \item \textbf{ViewModel (MVVM):} È l'approccio standard moderno. Il \texttt{ViewModel} "sopravvive" alle rotazioni. Lo stato della UI viene tenuto nel ViewModel.
    \item \textbf{\texttt{onSaveInstanceState()} e \texttt{onRestoreInstanceState()}:} Si salva un \texttt{Bundle} chiave-valore e lo si recupera alla ricreazione.
    \item \textbf{Forzare il Manifest (\texttt{android:configChanges}):} Specificando nel Manifest \texttt{android:configChanges="orientation|screenSize"}, Android \textit{non} riavvierà l'Activity, ma chiamerà \texttt{onConfigurationChanged()}. Questa pratica \textbf{non è consigliata} se non strettamente necessaria.
\end{enumerate}

\section{Tasks e Back Stack}
Un \textbf{Task} è un insieme di Activity con cui l'utente interagisce per svolgere un lavoro. Queste Activity sono organizzate in uno stack (LIFO) chiamato \textbf{Back Stack}.

Quando si avvia un'Activity con un Intent, si possono usare dei \textbf{Flags} per alterare il comportamento standard del Back Stack:
\begin{itemize}
    \item \texttt{Intent.FLAG\_ACTIVITY\_NEW\_TASK}: Avvia l'Activity in un nuovo task (o la porta in primo piano se il task esiste già). Indica che l'Activity funge da inizio di un nuovo lavoro.
    \item \texttt{Intent.FLAG\_ACTIVITY\_CLEAR\_TASK}: Causa la rimozione di tutte le Activity precedenti dal task corrente prima dell'avvio della nuova. Di solito usato in combinazione con \texttt{NEW\_TASK} (es. dopo un logout per impedire all'utente di tornare indietro).
\end{itemize}

\begin{center}
    \begin{tikzpicture}[
        node distance=0.5cm,
        box/.style={rectangle, draw, fill=blue!10, text width=2.5cm, align=center}
    ]
        \node[box] (a) {Activity A};
        \node[box, above=of a] (b) {Activity B};
        \node[box, above=of b, fill=green!20] (c) {Activity C (Top)};

        \draw[->, thick, red] (1.5, 1) -- node[right] {Tasto Back} (1.5, 0);
        \node[text width=5cm] at (5, 1) {Premendo "Back", Activity C viene distrutta (popped) e Activity B torna in stato di \texttt{onResume()}.};
    \end{tikzpicture}
\end{center}

\begin{tcolorbox}[colback=white,colframe=black,title=Comportamento dei Flag (Domande Quiz)]
    \begin{itemize}
        \item \textbf{\texttt{FLAG\_ACTIVITY\_NEW\_TASK}} e \textbf{\texttt{FLAG\_ACTIVITY\_CLEAR\_TASK}} usati insieme: Vengono rimosse tutte le Activity precedenti dal Task corrente e la nuova Activity diventa la root. Premendo "Back" si uscirà dall'app.
        \item Se l'utente preme il tasto Back: l'Activity in cima al Back Stack (foreground) viene distrutta (\texttt{onDestroy}) e la precedente torna in \texttt{onResume}.
    \end{itemize}
\end{tcolorbox}
    % --- START OF FILE 05_ui_liste.tex ---

\chapter{UI e Liste: Da ListView a Jetpack Compose}

La costruzione di interfacce grafiche in Android ha subito un'enorme evoluzione.
È vitale per l'esame conoscere sia l'approccio classico (XML e View) che quello moderno dichiarativo (Jetpack Compose).

\section{Approccio Classico: View e ViewGroup}
Il sistema classico si basa sull'albero degli oggetti UI, un'implementazione del \textbf{Design Pattern Composite}.
\begin{itemize}
    \item \textbf{\texttt{View}} (Component): La classe base per tutti gli elementi grafici interattivi (\texttt{TextView}, \texttt{Button}, ecc.). Implementa la funzione \texttt{onDraw()}.
    \item \textbf{\texttt{ViewGroup}} (Composite): Sottoclasse di \texttt{View} che funge da contenitore (Layout) per altre View o ViewGroup (es. \texttt{LinearLayout}, \texttt{FrameLayout}).
\end{itemize}

\subsection{Layout in XML}
Tra i \texttt{ViewGroup} più usati:
\begin{itemize}
    \item \textbf{\texttt{LinearLayout}:} Dispone gli elementi in una singola riga o colonna.
    Usa la proprietà \texttt{android:layout\_weight} per definire in proporzione \textit{quanto spazio extra} deve occupare un figlio rispetto agli altri.
    \item \textbf{\texttt{ConstraintLayout}:} Layout piatto (flat), basato su vincoli (constraint) tra gli elementi.
    Molto performante per UI complesse rispetto a LinearLayout innestati.
    \item \textbf{\texttt{FrameLayout}:} Progettato per contenere un singolo elemento principale.
    Se ha più figli, li \textbf{sovrappone} uno sopra l'altro (funziona da "cornice").
\end{itemize}

\section{Le Liste nel Sistema Classico: ListView e Adapters}
Per mostrare dati dinamici (es.
rubrica contatti) non si possono creare migliaia di TextView manualmente.
Si usano i componenti \texttt{AdapterView} (\texttt{ListView}, \texttt{GridView}) accompagnati dagli \textbf{Adapter}.

L'\textbf{Adapter} è una classe non-visiva che funge da "ponte" tra l'origine dei dati (Array, Database) e la \texttt{ListView}.
I principali sono \texttt{ArrayAdapter}, \texttt{CursorAdapter} e il generico \texttt{BaseAdapter}.

\subsection{Il concetto di View Recycling (Domanda da Esame!)}
Creare nuovi oggetti in Java/Kotlin (\textit{Object creation}) è un'operazione "costosa" in termini di memoria e CPU, e sollecita troppo il \textit{Garbage Collector}.
Per questo motivo, la \texttt{ListView} instanzia in RAM \textbf{solo il numero di righe strettamente necessarie} a riempire lo schermo, più un paio di scorta.

Quando l'utente fa scroll verso il basso:
\begin{enumerate}
    \item La riga che esce dallo schermo in alto finisce in una coda di riciclo (Scrap View).
    \item L'Adapter, nel suo metodo \texttt{getView(position, convertView, parent)}, riceve come parametro \texttt{convertView} la view appena uscita dallo schermo.
    \item Se \texttt{convertView != null}, significa che possiamo \textbf{riutilizzare le View esistenti}.
    Non facciamo un nuovo "inflate" (che consumerebbe CPU e RAM), ma ci limitiamo a cambiare i testi/immagini al suo interno. Questo riduce il consumo di memoria e migliora vertiginosamente le performance.
\end{enumerate}

\begin{center}
        % Scala il grafico per riempire la larghezza del testo (\textwidth)
    \resizebox{\textwidth}{!}{
        \begin{tikzpicture}[
            node distance=1.5cm,
            box/.style={rectangle, draw, rounded corners, fill=white, text width=4cm, align=center, drop shadow},
            arrow/.style={-Stealth, thick, blue}
        ]
            \node[box, fill=orange!10] (getview) {\textbf{Adapter.getView()}\\Chiesto elemento \#10};
            \node[box, right=of getview, fill=gray!10] (check) {\texttt{convertView == null}?};
            \node[box, top color=red!10, bottom color=red!20, above right=of check] (inflate) {\textbf{SI:} LayoutInflate\\(Costoso!)};
            \node[box, top color=green!10, bottom color=green!20, below right=of check] (reuse) {\textbf{NO:} Riutilizzo View\\(Riciclo, Molto veloce!)};
            \node[box, fill=blue!10, right=4cm of check] (bind) {Associazione Dati\\(Set text/image)};

            \draw[arrow] (getview) -- (check);
            \draw[arrow] (check) |- (inflate);
            \draw[arrow] (check) |- (reuse);
            \draw[arrow] (inflate) -| (bind);
            \draw[arrow] (reuse) -| (bind);
        \end{tikzpicture}
    }
\end{center}

\subsection{SimpleAdapter vs ArrayAdapter (Domanda Quiz)}
Mentre l'\texttt{ArrayAdapter} viene utilizzato principalmente quando i dati sono contenuti in una
\textbf{lista semplice di oggetti} o array, il \textbf{\texttt{SimpleAdapter}} si usa per dati più complessi.
\begin{itemize}
    \item \texttt{SimpleAdapter} è più flessibile perché \textbf{può associare più campi a diverse View} (es. un titolo, un sottotitolo e un'immagine nella stessa riga).
    \item Viene usato quando i dati sono \textbf{strutturati come Map (chiave-valore)}, tipicamente in una \texttt{List<Map<String, Any>>}.
    \item L'array di stringhe fornito nel suo costruttore rappresenta \textbf{le chiavi della Map}, mentre l'array di interi rappresenta \textbf{gli ID delle View nel layout} in cui inserire i dati.
    \item Se si tenta di caricare un'immagine da un URL web tramite SimpleAdapter, l'immagine \textbf{non viene caricata automaticamente}.
    È necessario implementare un \texttt{ViewBinder} personalizzato e usare librerie come Glide.
\end{itemize}

\subsection{L'evoluzione: RecyclerView}
Essendo la \texttt{ListView} limitata, Google ha introdotto la \textbf{\texttt{RecyclerView}}.
\begin{itemize}
    \item \textbf{Gestione del Layout:} Delega il posizionamento a un \texttt{LayoutManager} (es. \texttt{LinearLayoutManager} per liste, \texttt{GridLayoutManager} per griglie).
    \item \textbf{Performance:} Impone (forza) per architettura l'uso del pattern \textbf{ViewHolder}, obbligando lo sviluppatore a mantenere i riferimenti in memoria degli elementi della riga, rendendo il riciclo estremamente efficiente.
\end{itemize}

\section{Jetpack Compose: Il paradigma moderno}
Oggi, Google spinge \textbf{Jetpack Compose}, un toolkit UI \textbf{dichiarativo} interamente scritto in Kotlin.
Non ci sono più file XML, non c'è più `findViewById`, né si usano gli Adapter!

In Compose:
\begin{itemize}
    \item \textbf{Tutto è una funzione:} Le UI si creano scrivendo funzioni annotate con \texttt{@Composable}.
    I nomi di queste funzioni si scrivono in PascalCase (es. \texttt{fun MyButton()}).
    \item \textbf{UI come funzione dello Stato:} \texttt{UI = f(state)}.
    Se lo stato (i dati) cambia, Compose effettua automaticamente la \textbf{Ricomposizione} (\textit{Recomposition}), ridisegnando solo le parti di UI colpite dal cambiamento.
    \item \textbf{Gestione di liste lunghe:} Al posto di \texttt{RecyclerView}, si usano \texttt{LazyColumn} (lista verticale) e \texttt{LazyRow} (orizzontale).
    Queste funzioni creano componenti solo quando diventano visibili sullo schermo, garantendo le stesse performance della RecyclerView, ma in 3 righe di codice.
\end{itemize}

\begin{tcolorbox}[colback=yellow!10!white,colframe=yellow!70!black,title=Confronto XML vs Compose (Utile per i quiz)]
    \begin{itemize}
        \item \texttt{LinearLayout (vertical)} $\rightarrow$ \texttt{Column}
        \item \texttt{LinearLayout (horizontal)} $\rightarrow$ \texttt{Row}
        \item \texttt{FrameLayout} $\rightarrow$ \texttt{Box}
        \item \texttt{RecyclerView} $\rightarrow$ \texttt{LazyColumn} / \texttt{LazyRow}
        \item \texttt{ListView.setOnItemClickListener} $\rightarrow$ \texttt{Modifier.clickable \{ \}}
    \end{itemize}
\end{tcolorbox}

\begin{lstlisting}[caption={Creare una lista in Jetpack Compose}]
@Composable
fun ListaUtenti(utenti: List<String>) {
    // LazyColumn equivale a una RecyclerView efficiente
    LazyColumn {
        // Il costrutto 'items' sostituisce l'Adapter!
        items(utenti) { utente ->
            Text(
                text = utente,
                modifier = Modifier
                            .fillMaxWidth()
                            .padding(16.dp)
            )
        }
    }
}
\end{lstlisting}

\section{Caricamento di Immagini da Rete}
Sia con XML che con Compose, scaricare immagini da Internet sul \textit{Main Thread} farebbe bloccare o crashare l'app.
Nelle app in Compose moderne, al posto di gestire thread separati, si utilizza la libreria \textbf{Coil} (Coroutine Image Loader), usando direttamente il widget \texttt{AsyncImage}:
\begin{lstlisting}[caption={AsyncImage con Coil}]
AsyncImage(
    model = "https://picsum.photos/200/300",
    contentDescription = "Immagine di prova",
    modifier = Modifier.size(64.dp)
)
\end{lstlisting}
\texttt{AsyncImage} esegue il download in modo asincrono senza bloccare l'interfaccia, salva l'immagine in cache e la mostra appena pronta.
Ricorda che è sempre necessario il permesso \texttt{INTERNET} nell' \texttt{AndroidManifest.xml}.
    % --- START OF FILE 06_fragments_nav.tex ---

\chapter{Fragments e Navigation Component}

\section{Introduzione ai Fragment}
Un \textbf{Fragment} rappresenta \textbf{una porzione di interfaccia utente all'interno di un'Activity}.
Migliorano la riusabilità del codice e permettono di creare interfacce adattive per diversi dispositivi.
Nel ciclo di vita, il metodo utilizzato per creare la UI di un Fragment è \textbf{\texttt{onCreateView()}}, ed è al suo interno che il layout XML viene caricato tramite il metodo \textbf{\texttt{inflate()}}.

Per ospitare visivamente un Fragment in un layout XML, si utilizza il componente \textbf{\texttt{FragmentContainerView}}.

I Fragment offrono tre vantaggi principali:
\begin{itemize}
    \item \textbf{Modularità:} Dividono il codice complesso di un'Activity in parti più piccole e gestibili.
    \item \textbf{Riusabilità:} Permettono di riutilizzare la stessa porzione di UI (es.
    un menù laterale) in più Activity.
    \item \textbf{Adattabilità:} Consentono di mostrare layout diversi a seconda dell'orientamento o delle dimensioni dello schermo (es.
    Master-Detail flow su tablet).
\end{itemize}

\begin{tcolorbox}[colback=red!5!white,colframe=red!75!black,title=Nota per l'esame]
    Un Fragment \textbf{non può esistere da solo}.
    Deve essere \textbf{sempre} ospitato (embedded) all'interno di un'Activity.
    Il suo ciclo di vita è direttamente influenzato da quello dell'Activity host.
\end{tcolorbox}

\section{Il Ciclo di Vita del Fragment}
Il ciclo di vita di un Fragment è simile a quello di un'Activity, ma con alcuni metodi aggiuntivi cruciali.
Una tipica domanda d'esame riguarda la differenza tra la distruzione del Fragment e la distruzione della sua View.

\begin{center}
    \begin{tikzpicture}[
        node distance=0.6cm and 1cm,
        state/.style={rectangle, draw, rounded corners, fill=white, text width=4cm, align=center, drop shadow},
        arrow/.style={-Stealth, thick}
    ]
        \node[state, fill=gray!20] (attach) {\textbf{onAttach()}\\Assegnato all'Activity};
        \node[state, fill=blue!10, below=of attach] (create) {\textbf{onCreate()}\\Inizializzazione Fragment};
        \node[state, fill=green!10, below=of create, draw=green!50!black, thick] (createview) {\textbf{onCreateView()}\\Creazione gerarchia View};
        \node[state, fill=cyan!10, below=of createview] (activitycreated) {\textbf{onActivityCreated()}\\L'Activity host è pronta};
        \node[state, fill=yellow!20, below=of activitycreated] (start) {\textbf{onStart()} / \textbf{onResume()}\\Visibile e Attivo};

        \node[state, fill=orange!20, right=1.5cm of start] (pause) {\textbf{onPause()} / \textbf{onStop()}\\Non visibile};
        \node[state, fill=red!10, above=of pause, draw=red!50!black, thick] (destroyview) {\textbf{onDestroyView()}\\View distrutta (il Fragment vive)};
        \node[state, fill=red!20, above=of destroyview] (destroy) {\textbf{onDestroy()}\\Pulizia finale};
        \node[state, fill=gray!20, above=of destroy] (detach) {\textbf{onDetach()}\\Scollegato dall'Activity};

        \draw[arrow] (attach) -- (create);
        \draw[arrow] (create) -- (createview);
        \draw[arrow] (createview) -- (activitycreated);
        \draw[arrow] (activitycreated) -- (start);
        \draw[arrow] (start) -- (pause);
        \draw[arrow] (pause) -- (destroyview);
        \draw[arrow] (destroyview) -- (destroy);
        \draw[arrow] (destroy) -- (detach);

        \draw[arrow, dashed, blue] (destroyview.west) -- ++(-0.5,0) |- node[left, near start] {Ritorno dal Backstack} (createview.west);
    \end{tikzpicture}
\end{center}

\subsection{Metodi Fondamentali}
\begin{itemize}
    \item \textbf{\texttt{onCreateView()}:} È qui che si effettua l'\textit{inflate} del layout XML del fragment.
    Il metodo \textbf{deve ritornare la View} creata.
    \item \textbf{\texttt{onDestroyView()}:} La View associata al Fragment viene rimossa e distrutta. \textit{Attenzione:} Il Fragment stesso è ancora in memoria (potrebbe essere nel backstack)!
    \item \textbf{\texttt{onAttach()} / \textbf{onDetach()}:} Segnano il momento in cui il Fragment viene collegato o scollegato dall'Activity madre.
\end{itemize}

\begin{lstlisting}[caption={Creazione classica di un Fragment}, label={lst:fragment_create}]
class ExampleFragment : Fragment() {
    override fun onCreateView(
        inflater: LayoutInflater,
        container: ViewGroup?,
        savedInstanceState: Bundle?
    ): View? {
        // Inflate del layout XML per questo fragment
        return inflater.inflate(R.layout.fragment_example, container, false)
    }
}
\end{lstlisting}

\section{Gestione della Navigazione: Passato vs Presente}

\subsection{Il Vecchio Approccio: FragmentManager e FragmentTransaction}
Prima di Jetpack, cambiare Fragment richiedeva codice verboso e prono a errori.
Si utilizzava il \texttt{FragmentManager} per avviare una \texttt{FragmentTransaction}:
\begin{lstlisting}[caption={Vecchio approccio (Sconsigliato)}, label={lst:old_nav}]
val transaction = supportFragmentManager.beginTransaction()
transaction.replace(R.id.fragmentContainerView, DetailFragment())
transaction.addToBackStack(null) // Per gestire il tasto Indietro
transaction.commit()
\end{lstlisting}

\subsection{Il Nuovo Approccio: Navigation Component}
Oggi lo standard è l'uso del \textbf{Navigation Component} di Jetpack.
Elimina la necessità di scrivere transazioni manuali e gestisce automaticamente il \textit{Back Stack} (la pila delle schermate precedenti).

Si basa su 3 elementi chiave:
\begin{enumerate}
    \item \textbf{NavGraph (XML):} Un file XML (nella cartella \texttt{res/navigation}) in cui si definiscono tutte le destinazioni (Fragment) e le azioni (frecce di navigazione) tra di esse.
    \item \textbf{NavHost (\texttt{FragmentContainerView}):} Il contenitore vuoto all'interno del layout dell'Activity (\texttt{activity\_main.xml}) in cui i Fragment vengono scambiati.
    \item \textbf{NavController (Kotlin):} L'oggetto che nel codice gestisce effettivamente il cambio di schermata.
\end{enumerate}

\begin{lstlisting}[caption={Navigazione con NavController}, label={lst:new_nav}]
// Nel bottone del Fragment di partenza:
button.setOnClickListener {
    // Si naviga usando l'ID dell'azione o l'ID del fragment destinazione definito nel NavGraph
    findNavController().navigate(R.id.action_homeFragment_to_detailFragment)
}
\end{lstlisting}

\begin{tcolorbox}[colback=blue!5!white,colframe=blue!75!black,title=Domanda Tipica d'Esame]
    \textbf{Qual è il componente che si occupa di gestire materialmente lo swap dei fragment e il back stack nel Navigation Component?}\\
    \textit{Risposta:} Il \textbf{NavController}. Il \textit{NavHost} è solo il contenitore visivo, mentre il \textit{NavGraph} è solo la mappa (regole) della navigazione.
\end{tcolorbox}
    % --- START OF FILE 07_compose.tex ---

\chapter{Jetpack Compose: UI Moderna e Dichiarativa}

\section{Da XML a Jetpack Compose}
Storicamente, Android usava un approccio \textbf{Imperativo} per la UI: si creava un albero di View in XML e si modificava il loro stato tramite codice Kotlin/Java (usando \texttt{findViewById}, \texttt{setText}, ecc.).

\textbf{Jetpack Compose} sposa il paradigma \textbf{Dichiarativo}. L'interfaccia utente è scritta interamente in Kotlin tramite funzioni. Lo sviluppatore descrive \textit{cosa} la UI deve mostrare in base allo stato attuale, e Compose si occupa di aggiornarla automaticamente quando lo stato cambia.

\subsection{Regole base di Compose}
\begin{itemize}
    \item \textbf{Annotazione \texttt{@Composable}:} Qualsiasi funzione che emette UI deve essere annotata con \texttt{@Composable}.
    \item \textbf{PascalCase:} Per convenzione, le funzioni composable che fungono da widget si scrivono con l'iniziale maiuscola (es. \texttt{Greeting()}, \texttt{HomeScreen()}), comportandosi di fatto come se fossero delle classi grafiche.
    \item \textbf{Nessun XML:} Tutta la UI risiede in file \texttt{.kt}.
\end{itemize}

\section{Componenti Core e Layout}
I layout classici in XML hanno equivalenti diretti in Compose:

\begin{table}[h]
    \centering
    \begin{tabularx}{\textwidth}{|X|X|}
\hline
\rowcolor{gray!20} \textbf{Mondo XML (Imperativo)} & \textbf{Mondo Compose (Dichiarativo)} \\ \hline
\texttt{LinearLayout (vertical)} & \texttt{Column \{ \}} \\ \hline
\texttt{LinearLayout (horizontal)} & \texttt{Row \{ \}} \\ \hline
\texttt{FrameLayout} (Sovrapposizione) & \texttt{Box \{ \}} \\ \hline
\texttt{TextView} & \texttt{Text("Stringa")} \\ \hline
\texttt{RecyclerView} / \texttt{ListView} & \texttt{LazyColumn \{ \}} / \texttt{LazyRow \{ \}} \\ \hline
\texttt{setOnClickListener \{ \}} & \texttt{Modifier.clickable \{ \}} oppure \texttt{Button(onClick=...) \{ \}}\\ \hline
    \end{tabularx}
\end{table}

\subsection{I Modifier}
In Compose non esistono attributi XML come \texttt{android:layout\_width}. Tutto ciò che riguarda lo stile, le dimensioni, il padding e i comportamenti (click, scroll) viene passato tramite il parametro \textbf{\texttt{Modifier}}.

\begin{lstlisting}[caption={Esempio pratico di Modifier}, label={lst:modifier}]
Text(
    text = "Hello PDM!",
    modifier = Modifier
        .fillMaxWidth()        // Equivalente a match_parent
        .padding(16.dp)        // Margine/Padding
        .clickable { /* fai qualcosa */ }
)
\end{lstlisting}

\section{Lo Stato e la Ricomposizione (Cruciale per l'esame)}
In Compose, \textbf{la UI è una funzione dello Stato: UI = f(state)}.

\begin{center}
    \begin{tikzpicture}[
        node distance=2cm,
        box/.style={rectangle, draw, rounded corners, fill=white, text width=2.5cm, align=center, drop shadow},
        arrow/.style={-Stealth, thick}
    ]
        \node[box, fill=blue!10] (state) {\textbf{State}\\(es. variabile count)};
        \node[box, fill=green!10, below=of state] (ui) {\textbf{UI}\\(es. Text(count))};
        \node[box, fill=orange!10, right=of ui] (event) {\textbf{Event}\\(es. Button Click)};

        \draw[arrow, blue] (state) -- node[left] {Aggiorna} (ui);
        \draw[arrow, orange] (ui) -- node[below] {Genera} (event);
        \draw[arrow, red] (event) |- node[right] {Modifica} (state);
    \end{tikzpicture}
\end{center}

Se una semplice variabile Kotlin (es. \texttt{var count = 0}) cambia, la UI \textbf{non} si aggiorna. Per far capire a Compose che deve ri-disegnare la schermata, dobbiamo usare \textbf{\texttt{State}} e \textbf{\texttt{remember}}.

\begin{itemize}
    \item \textbf{\texttt{mutableStateOf(valore)}:} Crea un contenitore osservabile. Quando il valore cambia, Compose viene "notificato".
    \item \textbf{\texttt{remember \{ \}}:} Dice a Compose di "ricordare" questo valore tra una ricomposizione e l'altra, altrimenti ad ogni refresh la variabile verrebbe re-inizializzata al valore di partenza.
    \item \textbf{Ricomposizione (Recomposition):} È il processo con cui Compose ri-esegue le funzioni \texttt{@Composable} i cui dati di stato sono cambiati, aggiornando solo i pezzi di schermo necessari.
\end{itemize}

\begin{lstlisting}[caption={Gestione corretta dello Stato}, label={lst:compose_state}]
@Composable
fun Counter() {
    // 1. Dichiariamo lo stato
    var count by remember { mutableStateOf(0) }

    Button(onClick = {
        // 2. Modifichiamo lo stato (scatena la Ricomposizione)
        count++
    }) {
        // 3. La UI si aggiorna automaticamente
        Text("Click: $count")
    }
}
\end{lstlisting}

\section{Le Liste Dinamiche: LazyColumn}
Mentre nel vecchio sistema serviva creare una \texttt{RecyclerView}, un \texttt{Adapter}, un \texttt{ViewHolder} e un layout XML per la riga, in Compose le liste sono banali.

\textbf{\texttt{LazyColumn}} (e \texttt{LazyRow}) si chiamano "Lazy" perché instanziano (disegnano) solo gli elementi visibili sullo schermo in quel momento, offrendo grandissime performance.

\begin{lstlisting}[caption={Creare una lista in Compose}, label={lst:lazycolumn}]
@Composable
fun UserListScreen(users: List<String>) {
    LazyColumn {
        // Il blocco 'items' sostituisce il vecchio Adapter!
        items(users) { user ->
            // Qui definiamo la singola riga
            Text(text = "Utente: $user", modifier = Modifier.padding(8.dp))
        }
    }
}
\end{lstlisting}

\section{Esecuzione Asincrona e Immagini (Side Effects)}

\subsection{LaunchedEffect}
Se dobbiamo eseguire codice asincrono (come scaricare un JSON o dati da un database) non possiamo scriverlo direttamente nel corpo di un \texttt{@Composable}, altrimenti verrebbe rieseguito (scaricando ripetutamente i dati) ad ogni singola ricomposizione!

Per evitare questo si usa \textbf{\texttt{LaunchedEffect}}. Viene eseguito una sola volta quando il composable entra nello schermo (o quando una "key" specificata cambia). Per operazioni di rete si usa il blocco \texttt{withContext(Dispatchers.IO)}.

\subsection{Caricamento Immagini da URL}
In Jetpack Compose, per caricare un'immagine da Internet si consiglia la libreria di terze parti \textbf{Coil}, che offre il widget \textbf{\texttt{AsyncImage}}. Fa tutto da solo: apre un thread asincrono, scarica, salva in cache e mostra l'immagine, senza bloccare la UI.

\begin{lstlisting}[caption={Esempio di AsyncImage}, label={lst:asyncimage}]
AsyncImage(
    model = "https://picsum.photos/id/10/200/300",
    contentDescription = "Descrizione per accessibilita",
    modifier = Modifier.size(64.dp)
)
\end{lstlisting}

\section{Navigazione in Compose (Navigation Compose)}
Il Navigation Component è stato adattato per funzionare senza XML.
Servono 3 elementi:
\begin{enumerate}
    \item \textbf{\texttt{rememberNavController()}:} Per tenere traccia dello stato della navigazione.
    \item \textbf{\texttt{NavHost}:} Il contenitore che decide quale Composable mostrare. Sostituisce il NavGraph XML specificando una \texttt{startDestination}.
    \item \textbf{\texttt{composable("rotta")}:} Funzione che associa una stringa (la rotta) al Widget da disegnare.
\end{enumerate}

\begin{lstlisting}[caption={Navigazione in Compose}, label={lst:nav_compose}]
@Composable
fun AppNavigation() {
    val navController = rememberNavController()

    NavHost(navController = navController, startDestination = "home") {
        composable("home") {
            HomeScreen(onNavigate = { navController.navigate("detail") })
        }
        composable("detail") {
            DetailScreen()
        }
    }
}
\end{lstlisting}
    % --- START OF FILE 08_background_services.tex ---

\chapter{Esecuzione in Background: Services e Alternative}

Un'app Android ben progettata non si limita a reagire quando l'utente guarda lo schermo, ma deve saper gestire operazioni asincrone e di lunga durata in background (come riprodurre musica, scaricare file, sincronizzare dati).

\section{Cos'è un Service?}
Un \textbf{Service} è un componente dell'applicazione progettato per eseguire operazioni di lunga durata in background. \textbf{Non fornisce un'interfaccia utente (UI).}

\begin{tcolorbox}[colback=red!5!white,colframe=red!75!black,title=TRABOCCHETTO DA ESAME: Service vs Thread]
    Un errore comune è pensare che un Service giri in un thread separato. \textbf{Falso!}
    Di default, \textbf{un Service viene eseguito nel Main Thread (UI Thread)} dell'applicazione host. Se il Service deve eseguire operazioni bloccanti (es. chiamate di rete o calcoli intensivi), \textbf{devi creare tu un thread o una coroutine} all'interno del Service, altrimenti bloccherai l'interfaccia grafica scatenando un errore ANR (Application Not Responding).
\end{tcolorbox}

\section{Le tre tipologie di Service}
Storicamente, Android prevedeva 3 tipi di Service:
\begin{enumerate}
    \item \textbf{Background Service:} Esegue operazioni non direttamente visibili all'utente (es. comprimere lo storage). \textit{Attenzione:} Da Android 8.0 (API 26), il sistema impone forti limiti a questi service per risparmiare batteria, terminandoli forzatamente. Oggi si usa il \textbf{WorkManager}.
    \item \textbf{Foreground Service:} Esegue operazioni \textbf{evidenti all'utente} (es. player musicale o tracciamento GPS per la corsa). Per funzionare, l'app \textbf{deve mostrare una Notifica obbligatoria} nella barra di stato, in modo che l'utente sappia che l'app sta consumando risorse.
    \item \textbf{Bound Service:} Offre un'interfaccia Client-Server. Un'Activity (o un altro componente) si "lega" (bind) al Service per inviare richieste e ricevere risultati. Vive finché c'è almeno un client collegato ad esso.
\end{enumerate}

\section{Ciclo di vita: Started vs Bound}

Il ciclo di vita di un Service dipende da come viene avviato dall'Activity.

\begin{center}
    \begin{tikzpicture}[
        node distance=0.8cm and 2cm,
        state/.style={rectangle, draw, rounded corners, fill=white, text width=3cm, align=center, drop shadow},
        arrow/.style={-Stealth, thick}
    ]
        % Started Service (Left)
        \node[state, fill=blue!10] (startcall) {\texttt{startService()}};
        \node[state, below=of startcall] (oncreate1) {\texttt{onCreate()}};
        \node[state, fill=green!10, below=of oncreate1] (onstartcmd) {\texttt{onStartCommand()}\\(Il Service gira indefinitamente)};
        \node[state, fill=orange!10, below=of onstartcmd] (stopcall) {\texttt{stopService()} o \texttt{stopSelf()}};
        \node[state, fill=red!10, below=of stopcall] (ondestroy1) {\texttt{onDestroy()}};

        \draw[arrow] (startcall) -- (oncreate1);
        \draw[arrow] (oncreate1) -- (onstartcmd);
        \draw[arrow] (onstartcmd) -- (stopcall);
        \draw[arrow] (stopcall) -- (ondestroy1);

        % Bound Service (Right)
        \node[state, fill=blue!10, right=of startcall] (bindcall) {\texttt{bindService()}};
        \node[state, below=of bindcall] (oncreate2) {\texttt{onCreate()}};
        \node[state, fill=green!10, below=of oncreate2] (onbind) {\texttt{onBind()}\\(Ritorna IBinder)};
        \node[state, fill=orange!10, below=of onbind] (unbindcall) {\texttt{unbindService()}};
        \node[state, fill=red!10, below=of unbindcall] (ondestroy2) {\texttt{onDestroy()}};

        \draw[arrow] (bindcall) -- (oncreate2);
        \draw[arrow] (oncreate2) -- (onbind);
        \draw[arrow] (onbind) -- (unbindcall);
        \draw[arrow] (unbindcall) -- (ondestroy2);
    \end{tikzpicture}
\end{center}

\subsection{Implementare un Started Service (es. Foreground Music Player)}
Per implementare un Service che continua a suonare anche se l'app viene chiusa:
\begin{enumerate}
    \item Si dichiara il service nel \texttt{AndroidManifest.xml} con i permessi corretti (\texttt{FOREGROUND\_SERVICE} e il tipo specifico, es. \texttt{mediaPlayback}).
    \item Si estende la classe \texttt{Service}.
    \item Nel metodo \texttt{onStartCommand()} si lancia \texttt{startForeground(id, notification)} e si ritorna un intero che indica al sistema come comportarsi se il Service viene ucciso per mancanza di memoria:
    \begin{itemize}
        \item \textbf{\texttt{START\_STICKY}:} Il sistema \textbf{riavvia automaticamente} il Service appena possibile, passando un intent nullo (ideale per un music player).
        \item \textbf{\texttt{START\_NOT\_STICKY}:} Il Service \textbf{non} viene ricreato automaticamente, a meno che non ci siano nuovi intent in coda.
    \end{itemize}
\end{enumerate}

\begin{lstlisting}[caption={Creazione di un Foreground Service}]
class MusicForegroundService : Service() {
    override fun onStartCommand(intent: Intent?, flags: Int, startId: Int): Int {
        // Obbligatorio per i Foreground Service: mostrare una notifica
        startForeground(1, createNotification())

        // Avvio del player musicale (es. ExoPlayer/MediaPlayer)
        mediaPlayer?.start()

        // Se il sistema uccide il service per RAM, riavvialo appena possibile
        return START_STICKY
    }

    override fun onBind(intent: Intent?): IBinder? = null // Non e un Bound Service
}
\end{lstlisting}

\subsection{Implementare un Bound Service}
I Bound Service si usano quando l'Activity deve chiamare metodi specifici del Service (es. \texttt{musicService.pause()}). Per farlo, il Service restituisce un \textbf{\texttt{IBinder}} tramite il metodo \texttt{onBind()}. L'Activity usa un oggetto \texttt{ServiceConnection} per ricevere l'istanza del Binder.

\section{Modern Android: WorkManager vs Services}
A causa delle severe limitazioni introdotte nelle versioni moderne di Android (Doze Mode, App Standby), i classici \textit{Background Services} non si usano quasi più.

\textbf{Oggi la regola d'oro è:}
\begin{itemize}
    \item Task \textbf{immediato} ed evidente all'utente (Musica, Navigatore)? $\rightarrow$ \textbf{Foreground Service}.
    \item Task \textbf{differibile} o \textbf{periodico} (es. backup giornaliero, invio log, sync dati)? $\rightarrow$ \textbf{WorkManager}.
\end{itemize}

Il \textbf{WorkManager} è una libreria Jetpack che garantisce l'esecuzione del lavoro in background anche se l'app viene chiusa o il dispositivo viene riavviato. Rispetta le ottimizzazioni della batteria e permette di impostare dei \textit{Constraints} (es. "esegui solo se connesso in Wi-Fi e in ricarica").

\section{Altri Componenti di Background}
\begin{itemize}
    \item \textbf{Broadcast Receiver:} Componente senza UI che reagisce a messaggi di sistema (es. Wi-Fi connesso, batteria scarica). Registrabile via Manifest o da codice (\texttt{registerReceiver}).
    \item \textbf{Content Provider:} Struttura per condividere dati in modo sicuro tra app diverse. Usato tipicamente per accedere al database dei Contatti o della Galleria di sistema.
\end{itemize}
    % --- START OF FILE 09_persistenza.tex ---

\chapter{Persistenza dei Dati: Room e Firebase}

Un'app reale necessita di salvare dati in modo che sopravvivano alla chiusura o alla rotazione dello schermo.
In Android abbiamo diverse scelte:
\begin{itemize}
    \item \textbf{SharedPreferences / DataStore:} Per dati semplici chiave-valore (es.
    impostazioni, token di login).
    \item \textbf{File System:} Per dati non strutturati (es.
    foto, PDF).
    \item \textbf{Database (SQLite / Room):} Per dati strutturati, liste, relazioni e query complesse locali.
    \item \textbf{Cloud (Firebase):} Per dati condivisi tra più utenti/dispositivi.
\end{itemize}

\section{SQLite "Puro" (Il Motore)}
Android include nativamente \textbf{SQLite}, un database relazionale embedded (non c'è un processo server separato, ma una libreria C molto compatta e veloce).
Caratteristica peculiare di SQLite è il \textit{Type Affinity}: il controllo dei tipi è "rilassato" (puoi inserire una stringa in una colonna numerica).

L'approccio storico (ora considerato \textit{obsoleto} per la troppa verbosità) richiedeva:
\begin{enumerate}
    \item Estendere \texttt{SQLiteOpenHelper} per gestire creazione (\texttt{onCreate}) e aggiornamento schema (\texttt{onUpgrade}).
    \item Scrivere query SQL manuali, proni a errori a runtime (non validati dal compilatore).
    \item Usare \texttt{ContentValues} per gli \texttt{INSERT} e i \texttt{Cursor} (puntatori ai risultati) per leggere i dati estratti con query \texttt{SELECT}.
\end{enumerate}

\section{Jetpack Room: L'Approccio Moderno}
Per risolvere i problemi di SQLite puro (eccesso di codice \textit{boilerplate} e zero type-safety), Google ha introdotto \textbf{Room}, un layer di astrazione sopra SQLite.

Room è una \textbf{libreria che funziona come astrazione sopra SQLite}.
Si basa su 3 elementi fondamentali:
\begin{enumerate}
    \item \textbf{Entity:} Le data class Kotlin annotate con \texttt{@Entity}.
    In Room, una classe annotata con \texttt{@Entity} rappresenta \textbf{una tabella del database}.
    \item \textbf{DAO (Data Access Object):} Interfacce annotate con \texttt{@Dao}.
    In Room, il DAO è responsabile di \textbf{gestire le query}.
    I metodi annotati con \texttt{suspend} indicano che la funzione \textbf{deve essere chiamata in una coroutine}.
    Un metodo che ritorna un \texttt{Flow} serve a \textbf{emettere dati nel tempo} (aggiornando automaticamente la UI).
    \item \textbf{Database:} Una classe astratta annotata con \texttt{@Database}.
    Room controlla le query SQL \textbf{a compile-time}, garantendo sicurezza prima dell'esecuzione.
\end{enumerate}

\begin{center}
    \begin{tikzpicture}[
        node distance=0.8cm,
        box/.style={rectangle, draw, rounded corners, fill=white, text width=5cm, align=center, drop shadow},
        arrow/.style={-Stealth, thick}
    ]
        \node[box, fill=purple!10] (ui) {\textbf{UI (Jetpack Compose)}\\Osserva lo stato};
        \node[box, fill=blue!10, below=of ui] (vm) {\textbf{ViewModel}\\Raccoglie il Flow};
        \node[box, fill=green!10, below=of vm] (repo) {\textbf{Repository}\\Astrae la sorgente dati};
        \node[box, fill=orange!10, below=of repo] (dao) {\textbf{Room (DAO)}\\Converte oggetti in SQL};
        \node[box, fill=gray!20, below=of dao] (sqlite) {\textbf{SQLite}\\Motore fisico su disco};

        \draw[arrow] (ui) -- (vm);
        \draw[arrow] (vm) -- (repo);
        \draw[arrow] (repo) -- (dao);
        \draw[arrow] (dao) -- (sqlite);
    \end{tikzpicture}
\end{center}

\subsection{Configurazione e KSP}
Per far funzionare Room, Android deve "scrivere il codice SQLite al posto tuo" durante la compilazione.
Questo si fa usando il \textbf{KSP (Kotlin Symbol Processing)}.
Se nel file \texttt{build.gradle.kts} dimentichi di aggiungere il plugin KSP o la dipendenza \texttt{ksp(libs.androidx.room.compiler)}, il progetto non compilerà.

\subsection{Esempio di Codice Room}
\begin{lstlisting}[caption={I tre componenti di Room}]
// 1. ENTITY (Tabella)
@Entity(tableName = "shopping_items")
data class ShoppingItem(
    @PrimaryKey(autoGenerate = true) val id: Int = 0,
    val name: String,
    val quantity: Int
)

// 2. DAO (Query)
@Dao
interface ShoppingDao {
    // Flow trasforma la query in uno stream Reattivo!
    @Query("SELECT * FROM shopping_items")
    fun getAllItems(): Flow<List<ShoppingItem>>

    // suspend impone l'uso di Coroutine (no blocco UI)
    @Insert
    suspend fun insertItem(item: ShoppingItem)
}

// 3. DATABASE
@Database(entities = [ShoppingItem::class], version = 1)
abstract class AppDatabase : RoomDatabase() {
    abstract fun shoppingDao(): ShoppingDao
}
\end{lstlisting}

\begin{tcolorbox}[colback=blue!5!white,colframe=blue!75!black,title=Da Flow a UI Reattiva in Compose]
    Una potenza di Room + Compose è l'uso del \textbf{\texttt{Flow}}.
    Quando il DB viene aggiornato (es.
    inserimento), Room emette un nuovo \texttt{Flow}.
    Nella UI (Compose), usando la funzione \textbf{\texttt{collectAsStateWithLifecycle()}}, il \texttt{Flow} si converte in \texttt{State}.
    Al mutare dello State, Compose fa automaticamente la \textbf{Ricomposizione}, ridisegnando la lista senza bisogno di \texttt{notifyDataSetChanged()}!
\end{tcolorbox}

\section{Backend as a Service: Firebase}
Firebase è una piattaforma di Google che fornisce servizi di backend pronti all'uso.
Sostituisce la necessità di creare un proprio server, gestire database remoti e sviluppare API. Servizi principali: \textit{Authentication}, \textit{Cloud Storage} (per file), e i due database NoSQL.

\subsection{Realtime Database vs Cloud Firestore}
Una tematica classica da esame è la differenza tra i due database di Firebase.
Entrambi offrono sincronizzazione in tempo reale tramite \textbf{WebSocket}.

\begin{table}[h]
    \centering
    \begin{tabularx}{\textwidth}{|X|X|}
\hline
\rowcolor{gray!20} \textbf{Realtime Database} & \textbf{Cloud Firestore} \\ \hline
I dati sono memorizzati come un grande \textbf{JSON tree}. & I dati sono organizzati in \textbf{Collection} che contengono \textbf{Documenti}. Un document a sua volta può contenere \textbf{campi}. \\ \hline
Meno scalabile. Scaricare un nodo scarica \textit{tutti i dati} in esso contenuti (spesso si scaricano dati non necessari). & \textbf{Più scalabile}. Le query permettono il download mirato. Una query con più condizioni può richiedere la creazione di un \textbf{indice}. \\ \hline
Si usano listener come \texttt{ValueEventListener} per il tempo reale. & Il metodo \texttt{addSnapshotListener()} permette di ricevere aggiornamenti \textbf{in tempo reale}. \\ \hline
    \end{tabularx}
\end{table}

\subsection{Denormalizzazione (Best Practice per DB NoSQL)}
Nei database relazionali (SQL) si normalizzano i dati (es.
Join tables).
Nei database NoSQL come Firebase, la regola d'oro (per evitare di scaricare MegaBytes di messaggi solo per leggere il titolo
di una chat) è \textbf{appiattire i dati (Flattening) e Denormalizzare}.

Si evitano alberi profondi creando nodi separati al livello principale (top level nodes):
\begin{lstlisting}[language=SQL, caption={Antipattern vs Denormalizzazione}]
// SBAGLIATO (Nidificato)
"classes": {
    "class1": {
        "students": { "student1": true } // Se chiamo 'class1' scarico tutti gli studenti!
    }
}

// CORRETTO (Appiattito e Denormalizzato)
"classes": { "class1": {...dati classe} }
"students": { "student1": {...dati studente} }
"class_enrolments": {
    "class1": { "student1": true } // Tabella pivot usata solo per le query
}
\end{lstlisting}

\subsection{Lettura Dati e Listener}
Nei quiz è richiesto di riconoscere i listener appropriati:
\begin{itemize}
    \item \textbf{Realtime Database:} Si usa \texttt{ValueEventListener} (che ha i metodi \texttt{onDataChange} e \texttt{onCancelled}).
    \item \textbf{Firestore:} Per leggere una tantum si usa \texttt{.get().addOnSuccessListener}.
    Per ricevere aggiornamenti \textit{in tempo reale}, si usa \texttt{addSnapshotListener()}.
\end{itemize}

Le query in Firebase sono \textbf{Asincrone}.
Non bloccano la UI. I dati ottenuti vanno poi mappati in liste.
Se si usa il vecchio sistema XML con \texttt{SimpleAdapter}, bisogna ricordare che \textbf{non supporta nativamente il caricamento di immagini da URL}; per
farlo occorre definire un \texttt{ViewBinder} personalizzato e usare librerie come \textbf{Glide} o \textbf{Coil}.

\end{document}